\chapter{Management Approaches, SNMP and Syslog}\label{ch:management}
\bp{5.3,5.4,5.6}

\begin{outcomes}
By the end of this chapter you will be able to:
\begin{tight}
  \item \textbf{Describe} the five network management approaches the blueprint
        names, and tell controller-based from automation-based.
  \item \textbf{Describe} how SNMP works, what its three versions differ on, and
        when to use a trap rather than an inform.
  \item \textbf{Interpret} any IOS syslog message: facility, severity, mnemonic
        and description.
  \item \textbf{Recite} the eight severity levels and explain what a logging
        threshold hides.
\end{tight}
\end{outcomes}

\begin{prereq}
\chref{aaa} for secure management access and for why credentials in clear are
unacceptable. \chref{transport} for UDP. \chref{tshootl2} --- you have already
read syslog output there; this chapter makes it systematic.
\end{prereq}

\begin{keyterms}
device-based $\bullet$ controller-based $\bullet$ cloud-based $\bullet$
automation-based $\bullet$ infrastructure as code $\bullet$ SNMP manager
$\bullet$ agent $\bullet$ MIB $\bullet$ OID $\bullet$ community string $\bullet$
trap $\bullet$ inform $\bullet$ facility $\bullet$ severity $\bullet$ mnemonic
\end{keyterms}

\begin{examnote}[Three topics, three different verbs]
\tp{5.3} \emph{describe network management approaches (device-based,
cloud-based, controller-based, automation-based, and infrastructure as code)} ---
five named approaches, so expect a matching or scenario question.

\tp{5.4} \emph{describe the function of SNMP in network operations} --- function,
not configuration. What it does and how, not how to type it.

\tp{5.6} \emph{interpret syslog message content, severity levels, and facilities}
--- \textbf{interpret}. You will be shown a log line and asked what it means.
Section~\ref{sec:syslog-read} is built for exactly that and is the highest-value part of this
chapter.
\end{examnote}

\section{Five ways to manage a network}

\begin{longtable}{@{}L{3.1cm}L{5.7cm}L{6.1cm}@{}}
\toprule
\rowcolor{primary!15}
\textbf{Approach} & \textbf{How it works} & \textbf{Trade-off} \\
\midrule
\endhead
\textbf{Device-based} &
You connect to each device individually by console or SSH and configure it &
Complete control and no dependencies. Does not scale, and every device is a
separate opportunity to make a different mistake \\
\rowcolor{lightbg}
\textbf{Controller-based} &
A controller holds the intended state and programs the devices to match
(Catalyst Center, SD-WAN, SD-Access) &
Consistency and central policy. The controller becomes critical
infrastructure, and you manage the network its way \\
\textbf{Cloud-based} &
The management plane is a service the vendor runs; devices call home to it
(Meraki) &
Nothing to host, upgrade or back up. You depend on the vendor's availability
and on internet reachability from every site \\
\rowcolor{lightbg}
\textbf{Automation-based} &
Scripts and tools drive the devices' existing interfaces --- Ansible, Python
(\chref{ansible}) &
Repeatable, and adds no new control plane. The devices are unchanged, so it
works with what you already own \\
\textbf{Infrastructure as code} &
The desired configuration lives in version-controlled files and is applied by a
pipeline &
Reviewable, testable, revertible; the repository is the source of truth.
Requires discipline and tooling a small team may not have \\
\bottomrule
\end{longtable}

\begin{alertbox}[Controller-based against automation-based: the distinction that gets tested]
They sound similar and are structurally different.

\textbf{Controller-based} inserts a \emph{new system} between you and the
network. The controller holds the intended state, and it is in the path
permanently. Take it away and you lose central management.

\textbf{Automation-based} adds \emph{tooling} outside the network. Ansible logs
in over SSH and types commands faster and more consistently than you can. Take it
away and every device still works exactly as before, because nothing about them
changed.

One-line test: \textbf{does the network still have a management system when the
tool is switched off?} Controller, yes --- and you have lost it. Automation, no
--- and you have lost nothing but convenience.

Infrastructure as code is a step beyond automation: not just \emph{running} the
change, but keeping the intended state in a repository that is reviewed, tested
and versioned.
\end{alertbox}

\section{SNMP}

\begin{definitionbox}[SNMP]
The Simple Network Management Protocol. A \term{manager} polls \term{agents}
running on network devices for values held in a \term{MIB} --- a tree of
variables, each identified by a numeric \term{OID}. Agents listen on UDP~161;
the manager receives unsolicited notifications on UDP~162.
\end{definitionbox}

\ccnafig{%
\begin{tikzpicture}[font=\scriptsize]
  \tikzset{
    dv/.style={rounded corners=3pt, draw=#1, fill=#1!8, text=#1!45!black,
               line width=0.9pt, minimum width=1.7cm, minimum height=1.1cm,
               align=center, font=\scriptsize\bfseries},
    a/.style={-{Stealth[length=2.4mm]}, line width=1pt}}

  \node[dv=dom5] (nms) at (0,0) {\faIcon{desktop}\\NMS\\\tiny\mdseries manager};
  \node[dv=dom2] (sw)  at (7.6,1.2) {\faIcon{sitemap}\\SW1\\\tiny\mdseries agent};
  \node[dv=dom3] (r)   at (7.6,-1.2) {\faIcon{route}\\R1\\\tiny\mdseries agent};

  \draw[a, dom1!75] (nms) -- node[font=\tiny, above, midway, sloped,
        text=dom1!45!black] {Get / GetNext / GetBulk $\rightarrow$ UDP 161} (sw);
  \draw[a, dom4!75] (r) -- node[font=\tiny, below, midway, sloped,
        text=dom4!45!black] {Trap / Inform $\rightarrow$ UDP 162} (nms);

  % Below the diagram, not between the arrows. Centred at (3.8,0.15) this
  % legend is seven centimetres wide, so it started inside the NMS node and
  % ran straight through both arrows -- "polling" was struck through by the
  % line it describes.
  \node[font=\tiny, text=neutral!85!black, align=center] at (3.8,-2.35)
      {\textbf{polling} --- the manager asks, on a schedule\\[2pt]
       \textbf{notification} --- the agent speaks when something happens};
\end{tikzpicture}}%
{SNMP works in both directions: the manager polls on port 161, and agents report
events unprompted to port 162.}{snmp}

\subsection{The operations}

\begin{tight}
  \item \textbf{Get} --- fetch one value. \textbf{GetNext} --- fetch the next one
        in the tree, which is how a manager walks a whole table.
  \item \textbf{GetBulk} --- fetch many at once. Added in v2c, and the reason
        polling a large switch is not painfully slow.
  \item \textbf{Set} --- \emph{write} a value, changing the device's
        configuration or state. Powerful, and the reason read-write access needs
        treating as an administrative credential.
  \item \textbf{Trap} --- the agent notifies the manager of an event. Fire and
        forget: sent once, over UDP, never acknowledged.
  \item \textbf{Inform} --- the same, but \emph{acknowledged} and retransmitted
        until it is. Added in v2c.
\end{tight}

\begin{conceptbox}[Trap or inform]
A trap that is lost is lost silently, and the event it described never reaches
anyone. An inform costs a round trip and some agent memory, and survives a
dropped packet.

Use informs for anything you would be sorry to miss --- a device losing power
redundancy, a security event. Use traps for high-volume, low-consequence
notifications where losing one does not matter. This is a fair exam question and
a real design decision.
\end{conceptbox}

\subsection{The three versions}

\begin{longtable}{@{}L{2.0cm}L{4.6cm}L{4.0cm}L{4.3cm}@{}}
\toprule
\rowcolor{primary!15}
\textbf{Version} & \textbf{Authentication} & \textbf{Encryption} & \textbf{Verdict} \\
\midrule
\endhead
v1 & Community string, sent in clear &
None & Obsolete. No GetBulk, no informs \\
\rowcolor{lightbg}
v2c & Community string, sent in clear &
\textbf{None} & Still widely deployed. Adds GetBulk and informs. The community
string is a password on the wire in plain text \\
v3 & Username with HMAC-SHA authentication &
\textbf{AES} & The only version to use. Three levels: noAuthNoPriv, authNoPriv,
authPriv \\
\bottomrule
\end{longtable}

\begin{alertbox}[``v2c is secure because the community string is not \cmd{public}'']
It is not. The community string travels in clear text in every single polling
packet, hundreds of times an hour, and anyone who can capture one packet has it.
Renaming it from \cmd{public} to something obscure changes nothing.

And a read-write community string is equivalent to an enable password: an
attacker holding one can reconfigure the device, and on many platforms can pull
the entire running configuration. If SNMP must be v2c for a legacy tool, make it
read-only, restrict it with an ACL to the management station's address, and plan
its removal.
\end{alertbox}

\begin{config}{SNMPv3, read-only, for a monitoring station}
snmp-server view OPSVIEW iso included
snmp-server group OPSGROUP v3 priv read OPSVIEW
!
! auth protects integrity and identity; priv adds encryption.
! "priv" implies "auth" -- you cannot encrypt without authenticating.
snmp-server user nms OPSGROUP v3 auth sha AuthPassPhrase priv aes 128 PrivPassPhrase
!
snmp-server host 10.0.0.45 version 3 priv nms
snmp-server enable traps
!
! These two are free and save an hour when somebody has to find the
! device in a building.
snmp-server location Campus Block C, Rack 4, RU 22
snmp-server contact netops@example.edu
\end{config}

\begin{conceptbox}[What SNMP is actually used for]
Three jobs, and the exam's phrase ``in network operations'' points at them.

\begin{tight}
  \item \textbf{Performance monitoring} --- polling interface counters every few
        minutes to produce the traffic graphs everybody recognises.
  \item \textbf{Fault notification} --- traps and informs for events that cannot
        wait for the next poll.
  \item \textbf{Inventory} --- serial numbers, models, software versions and
        uptime collected across the estate.
\end{tight}

What it is \emph{not} good at is configuration. \cmd{Set} exists, but nobody
manages configuration by SNMP; that is what Chapters 31 and the automation
approaches are for.
\end{conceptbox}

\section{Syslog, and how to read a message}\label{sec:syslog-read}

Syslog is where a device tells you what it did. Reading it fluently is the single
most transferable skill in this chapter.

\subsection{Anatomy of an IOS message}

\ccnafig{%
\begin{tikzpicture}[font=\scriptsize]
  % Each part of the message is its own node, chained left to right, so
  % the braces below can attach to real anchors instead of to guessed
  % character positions.
  \tikzset{
    seg/.style={font=\ttfamily\small, inner xsep=0pt, inner ysep=2pt,
                anchor=base west},
    br/.style={decorate, decoration={brace, amplitude=3pt, mirror},
               line width=0.7pt, #1},
    segcap/.style={font=\tiny\bfseries, text=#1!45!black, align=center,
                anchor=north}}

  \node[seg]                        (ts)  at (0,0)             {Aug 27 09:14:22.331:~};
  \node[seg, text=dom1!45!black]    (fac) at (ts.base east)    {\%LINEPROTO};
  \node[seg]                        (d1)  at (fac.base east)   {-};
  \node[seg, text=accent!70!black]  (sev) at (d1.base east)    {5};
  \node[seg]                        (d2)  at (sev.base east)   {-};
  \node[seg, text=dom3!45!black]    (mne) at (d2.base east)    {UPDOWN};
  \node[seg]                        (cl)  at (mne.base east)   {:};

  % -0.30cm put this line's top inside the space the brace labels below the
  % first row occupy, so "timestamp", "facility" and "mnemonic" printed on
  % top of the description and both were unreadable. The labels hang about
  % 13pt below seg.south; leave clear air under that.
  \node[seg, text=dom2!45!black, anchor=north west] (desc)
      at ([yshift=-0.68cm]ts.south west)
      {Line protocol on Interface GigabitEthernet0/1, changed state to down};

  \draw[br=neutral!70] ([yshift=-2pt]ts.south west) -- ([yshift=-2pt]ts.south east);
  \node[segcap=neutral] at ([yshift=-6pt]ts.south) {timestamp};

  \draw[br=dom1!80] ([yshift=-2pt]fac.south west) -- ([yshift=-2pt]fac.south east);
  \node[segcap=dom1] at ([yshift=-6pt]fac.south) {facility};

  \draw[br=dom3!80] ([yshift=-2pt]mne.south west) -- ([yshift=-2pt]mne.south east);
  \node[segcap=dom3] at ([yshift=-6pt]mne.south) {mnemonic};

  % Severity is one character wide, so it gets a leader upward instead
  % of a brace it would dwarf.
  \draw[accent!80, line width=0.7pt] (sev.north) -- ++(0,0.34)
      node[font=\tiny\bfseries, text=accent!70!black, anchor=south] {severity};

  \draw[br=dom2!80] ([yshift=-2pt]desc.south west) -- ([yshift=-2pt]desc.south east);
  \node[segcap=dom2] at ([yshift=-6pt]desc.south) {description --- the part in English};
\end{tikzpicture}}%
{Every IOS syslog message has the same four parts after the timestamp. The
severity is the single digit between the two hyphens.}{syslogfmt}

\begin{alertbox}[The severity is the digit in the middle, and it is not the number of anything else]
In \cmd{\%LINEPROTO-5-UPDOWN} the \cmd{5} is the severity level. In
\cmd{\%LINK-3-UPDOWN} the \cmd{3} is the severity level.

Those two messages describe the same interface going down and carry different
severities because they concern different layers. Seeing both together is normal
and tells you the physical link dropped, which took the line protocol with it. A
\cmd{\%LINEPROTO-5-UPDOWN} \emph{without} a matching \cmd{\%LINK-3-UPDOWN} means
layer 1 stayed up and something above it failed --- keepalives, an encapsulation
mismatch, or a shutdown of a subinterface.
\end{alertbox}

\subsection{The eight severity levels}

\begin{longtable}{@{}L{1.3cm}L{2.9cm}L{5.3cm}L{5.4cm}@{}}
\toprule
\rowcolor{primary!15}
\textbf{Level} & \textbf{Name} & \textbf{Meaning} & \textbf{Typical IOS message} \\
\midrule
\endhead
\rowcolor{accent!12}
\textbf{0} & Emergency & System unusable & Imminent thermal shutdown \\
\rowcolor{accent!8}
\textbf{1} & Alert & Immediate action required & Redundant power supply failed \\
\rowcolor{accent!5}
\textbf{2} & Critical & Critical condition &
\cmd{\%SPANTREE-2-ROOTGUARD\_BLOCK} \\
\textbf{3} & Error & Error condition &
\cmd{\%LINK-3-UPDOWN} --- interface down \\
\rowcolor{lightbg}
\textbf{4} & Warning & Warning condition &
\cmd{\%PM-4-ERR\_DISABLE} --- port err-disabled \\
\textbf{5} & Notification & Normal but significant &
\cmd{\%SYS-5-CONFIG\_I} --- somebody saved a change \\
\rowcolor{lightbg}
\textbf{6} & Informational & Informational message &
\cmd{\%SEC-6-IPACCESSLOGP} --- an ACL matched \\
\textbf{7} & Debugging & Debug output &
Anything from a \cmd{debug} command \\
\bottomrule
\end{longtable}

\begin{conceptbox}[Remembering the order]
\textbf{E}very \textbf{A}wesome \textbf{C}isco \textbf{E}ngineer \textbf{W}ill
\textbf{N}eed \textbf{I}ce-cream \textbf{D}aily --- Emergency, Alert, Critical,
Error, Warning, Notification, Informational, Debugging.

Two facts that follow and are examined: \textbf{0 is the most severe and 7 the
least}, which is the opposite of most people's instinct. And configuring a
logging destination at level \emph{n} sends everything from 0 \emph{through}
\emph{n} --- so \cmd{logging console warnings} (level 4) also delivers levels 0
to 3, and silently discards 5, 6 and 7.
\end{conceptbox}

\subsection{Configuring logging}

\begin{config}{a sensible logging setup}
! Without this every message is stamped only with uptime, which is
! useless for correlating across devices. Do this first, everywhere.
service timestamps log datetime msec localtime show-timezone
service timestamps debug datetime msec localtime show-timezone
!
! Keep a generous local buffer at informational, so the evidence is
! on the device even if the syslog server was unreachable.
logging buffered 65536 informational
!
! Console logging on a busy device makes the CLI unusable. Turn it
! down; the buffer and the server still have everything.
logging console warnings
logging monitor informational
!
! The central server. "trap" here means the syslog severity
! threshold, NOT an SNMP trap -- an unfortunate collision of terms.
logging host 10.0.0.45
logging trap notifications
logging source-interface Loopback0
!
line con 0
 logging synchronous
line vty 0 15
 logging synchronous
\end{config}

\begin{alertbox}[\cmd{logging trap} has nothing to do with SNMP traps]
It sets the \emph{severity threshold for messages sent to a syslog server}.
\cmd{logging trap notifications} means ``send levels 0 to 5 to the host''. The
word is a historical accident and it confuses people every year.

\cmd{logging synchronous} is unrelated to both: it stops a log message printed
mid-command from scrambling what you were typing. It changes nothing about what
is logged and it will save your temper.
\end{alertbox}

\begin{verify}{SW1}{show logging | include Console|Buffer|Trap}
    Console logging: level warnings, 2201 messages logged
    Monitor logging: level informational, 0 messages logged
    Buffer logging:  level informational, 88214 messages logged
    Trap logging: level notifications, 41022 message lines logged
\end{verify}

\begin{pitfall}
\begin{tight}
  \item \textbf{Reading severity backwards.} 0 is worst, 7 is least important.
  \item \textbf{Setting a threshold too high and concluding nothing is wrong.}
        The message may exist and be filtered. Check the levels before trusting
        an empty log.
  \item \textbf{No \cmd{service timestamps}.} Messages stamped with uptime cannot
        be correlated across devices or with a user's report.
  \item \textbf{Confusing \cmd{logging trap} with SNMP traps.}
  \item \textbf{Leaving console logging at debugging on a production device.} The
        CLI becomes unusable exactly when you need it.
  \item \textbf{SNMP v2c with a read-write community string.} That is an enable
        password sent in clear thousands of times a day.
  \item \textbf{No \cmd{logging source-interface}.} Messages arrive from whichever
        interface routing chose, and the server cannot reliably attribute them.
  \item \textbf{Relying on traps for events that matter.} Use informs.
  \item \textbf{Calling Ansible ``controller-based''.} It adds no control plane.
\end{tight}
\end{pitfall}

\begin{breakfix}[Ports keep going down overnight. The log ``shows nothing''.]
An access switch drops several user ports most nights. By morning they are back.
The site's monitoring shows the interfaces bouncing, but the network team reports
that the device's own log contains no explanation.

\begin{verify}{SW1}{show logging | include level}
    Console logging: level errors, 14 messages logged
    Monitor logging: level errors, 0 messages logged
    Buffer logging:  level errors, 14 messages logged
    Trap logging: level errors, 14 message lines logged
\end{verify}

\textbf{Diagnosis.} Every destination is set to \cmd{errors}, which is severity~3.
That delivers levels 0 to 3 and \textbf{discards 4, 5, 6 and 7}.

Look at what lives in the discarded range:

\begin{tight}
  \item \cmd{\%PM-4-ERR\_DISABLE} --- severity 4. The port being disabled.
  \item \cmd{\%LINEPROTO-5-UPDOWN} --- severity 5. The line protocol dropping.
  \item \cmd{\%SYS-5-CONFIG\_I} --- severity 5. Somebody making a change.
\end{tight}

Almost everything that explains why a port went down is severity 4 or 5, and none
of it is being recorded anywhere --- not the buffer, not the console, not the
server.

The log does not ``show nothing''. The log was told not to look.

The 14 messages that did survive are the severity-3 \cmd{\%LINK-3-UPDOWN} events
--- the interfaces going down. So the team can see the effect and has filtered out
every possible cause, which is precisely the worst configuration to have.

\textbf{Fix.}

\begin{config}{recording enough to diagnose}
service timestamps log datetime msec localtime show-timezone
logging buffered 65536 informational
logging trap notifications
logging console warnings
\end{config}

Then wait for the next occurrence rather than guessing. With level~5 reaching the
buffer and the server, the sequence around each drop becomes visible --- and in
this case it showed \cmd{\%PM-4-ERR\_DISABLE} for \cmd{psecure-violation}
following a cleaner's floor polisher being plugged into a wall port each night.

\textbf{The lesson.} \textbf{An empty log is a claim about the logging
configuration, not about the network.} Before concluding that nothing happened,
run \cmd{show logging} and read the levels. This is the most useful habit in the
whole chapter.
\end{breakfix}

\begin{lab}{Making a device tell you what it is doing}{Packet Tracer for syslog; CML for SNMPv3}
\textbf{Platform note.} Packet Tracer has a syslog server and will carry tasks
1--6. SNMPv3 needs a real image and a real NMS, so tasks 7--9 want CML or GNS3
with an SNMP tool on a Linux host.

\begin{tasks}
  \item On a switch with default logging, run \cmd{show logging} and record every
        threshold before changing anything.
  \item Enable \cmd{service timestamps} with date, time and milliseconds. Shut an
        interface and compare the message with one logged before the change.
  \item Shut and unshut an interface and capture both messages. Identify the
        facility, severity and mnemonic in each, and explain why the two
        severities differ.
  \item Set \cmd{logging buffered errors}. Repeat the shut and unshut. Record
        which of the two messages survives and explain why --- this is the Break
        \& Fix in miniature.
  \item Restore \cmd{logging buffered informational}, save a configuration
        change, and find the \cmd{\%SYS-5-CONFIG\_I} message. Note that it records
        \emph{that} a change was made, not what it was.
  \item Point \cmd{logging host} at a syslog server, set \cmd{logging trap
        notifications}, and confirm messages arrive. Then set it to
        \cmd{emergencies} and confirm they stop.
  \item \textbf{SNMP.} Configure SNMPv3 with authPriv and poll the device's
        interface table from a Linux host. Record the OID you used.
  \item Capture the polling traffic. Confirm you cannot read it. Then reconfigure
        for v2c, capture again, and \textbf{find the community string in the
        packet}. Write down what you saw.
  \item Configure a trap receiver, generate an event, and confirm receipt. Change
        the trap to an inform and observe the acknowledgement.
  \item \textbf{Extension.} Write the logging and SNMP standard your site should
        apply to every device, with a one-line justification per command, and the
        one thing you would check first when someone says the log is empty.
\end{tasks}
\end{lab}

\begin{checkpoint}
\begin{tightnum}
  \item Name the five management approaches and give the one-line test that
        separates controller-based from automation-based.
  \item In \cmd{\%PM-4-ERR\_DISABLE}, identify the facility, the severity and the
        mnemonic, and say what the severity means.
  \item Which severity levels does \cmd{logging buffered warnings} record, and
        which does it discard?
  \item State the difference between an SNMP trap and an inform, and when you
        would choose each.
  \item Why is SNMP v2c unsuitable regardless of how obscure the community string
        is?
  \item What does \cmd{logging trap} configure, and what does it not?
\end{tightnum}
\end{checkpoint}

\begin{chaptersummary}
\begin{tight}
  \item Five approaches: device-based (per box), controller-based (a controller
        holds intent), cloud-based (vendor-hosted management plane),
        automation-based (tooling drives existing interfaces), infrastructure as
        code (state in a reviewed repository).
  \item Controller-based adds a system to the network; automation-based does not.
        That is the distinction to be able to state.
  \item SNMP: manager polls agents on UDP~161 with Get, GetNext, GetBulk and Set;
        agents send traps and informs to UDP~162. Informs are acknowledged, traps
        are not.
  \item v1 and v2c send community strings in clear. Only v3 authenticates and
        encrypts. Use authPriv.
  \item An IOS syslog message is
        \cmd{\%FACILITY-SEVERITY-MNEMONIC: description}. The severity is the
        digit between the hyphens.
  \item Levels 0 to 7, emergency to debugging. \textbf{0 is worst.} A threshold
        of \emph{n} records 0 through \emph{n} and discards the rest.
  \item \cmd{service timestamps} first, always. \cmd{logging trap} is a syslog
        severity, not SNMP.
  \item An empty log describes the logging configuration, not the network.
\end{tight}
\end{chaptersummary}

\begin{reviewq}
  \item \textbf{(MCQ)} In \cmd{\%LINEPROTO-5-UPDOWN}, what does the \cmd{5}
        indicate?
        \begin{tight}
          \item[(a)] The interface number
          \item[(b)] The severity level, notification
          \item[(c)] The number of times the event occurred
          \item[(d)] The facility code
        \end{tight}
  \item \textbf{(MCQ)} Which severity level is the most severe?
        \begin{tight}
          \item[(a)] 7, debugging \quad (b)~0, emergency
          \quad (c)~1, alert \quad (d)~4, warning
        \end{tight}
  \item \textbf{(MCQ)} Which SNMP notification is acknowledged by the receiver?
        \begin{tight}
          \item[(a)] Trap \quad (b)~Inform \quad (c)~GetBulk \quad (d)~Set
        \end{tight}
  \item \textbf{(Short)} Explain the difference between controller-based and
        automation-based management, and give one example of each.
  \item \textbf{(Short)} A colleague reports that a switch's log contains no
        record of an outage that monitoring clearly shows. Give the first command
        you would run and the two things it might reveal.
  \item \textbf{(Applied)} Interpret this line fully, naming every component and
        explaining what happened and what you would do next:
        \cmd{Aug 27 02:41:09.882: \%PM-4-ERR\_DISABLE: psecure-violation error
        detected on Gi1/0/14, putting Gi1/0/14 in err-disable state}
  \item \textbf{(Applied)} Write the logging configuration for a production access
        switch that keeps full evidence locally, sends notifications and above to
        a central server, and does not make the console unusable. Justify each
        threshold you choose.
\end{reviewq}
